\phantomsection
\addcontentsline{toc}{chapter}{Abstract}
\begin{center}
\textbf{\large{Abstract}}
\end{center}

\begin{small}
The thesis \textit{Research on Intrinsic Motivation for Reinforcement Learning} studies reinforcement learning in sparse-reward environments, where exploration is difficult because feedback from the environment is rare. Many recent methods augment the extrinsic reward with curiosity- or novelty-based intrinsic rewards, but typically combine the two through a fixed coefficient. A fixed coefficient is sensitive to the specific task and ignores the fact that exploration needs differ across states.

The thesis proposes \textbf{Correlation-Aware Reward Exploration (CARE)}, a state-dependent intrinsic-reward scaling mechanism realized through a lightweight Beta Network. The Beta Network is trained on a correlation objective that maximizes the within-batch correlation between the weighted intrinsic reward and the extrinsic GAE advantage of the value function. Intrinsic reward is therefore amplified at states where the intrinsic signal co-varies with positive extrinsic advantage---states where the policy is still outperforming its own value baseline---and attenuated elsewhere. The method is applied to three representative intrinsic-reward families: count-based exploration, the Intrinsic Curiosity Module (ICM), and Rewarding Impact-Driven Exploration (RIDE), all combined with Proximal Policy Optimization (PPO) as the base policy optimizer.

Experimental evaluation on six sparse-reward MiniGrid environments shows that CARE matches or exceeds the best post-hoc fixed coefficient on most count-based and RIDE pairings without environment-specific tuning, while inheriting alignment issues on ICM where the underlying signal does not co-vary with extrinsic learning slack. When the extrinsic GAE advantage carries no signal-strength information across states, the correlation gradient vanishes and the scaling factor falls back to its prior $\beta_0$, reducing CARE to fixed-coefficient behaviour.

\vspace*{1cm}
\textbf{Keywords}: reinforcement learning, sparse rewards, intrinsic motivation, curiosity-driven exploration, adaptive reward scaling, ICM, RIDE
\end{small}
